\section{模型评估}

\subsection{评估的重要性}

讲者坦言 LLM 评估是整个领域最大的难题之一——"我们其实不太知道自己在做什么"。但评估依然至关重要，因为 Post-Training 的本质是\textbf{优化}，如果评估指标不对，优化方向就会偏离。

\subsection{评估方法一：Automated Benchmarks}

\textbf{原理}：给定问题集（如 MMLU 的多选题），模型选择答案，以准确率等指标计分。

\begin{table}[H]
\centering
\caption{常见自动化 Benchmark}
\begin{tabular}{lll}
\toprule
\textbf{Benchmark} & \textbf{领域} & \textbf{指标} \\
\midrule
MMLU & 通用知识（57个学科） & 准确率 \\
GSM8K & 小学数学推理 & 准确率 \\
MATH & 竞赛数学 & 准确率 \\
HumanEval & 代码生成 & pass@k \\
IFEval & Instruction Following & 约束满足率 \\
\bottomrule
\end{tabular}
\end{table}

\textbf{优点}：可扩展、低成本、可复现、可针对特定任务设计。

\textbf{缺点}：不反映模型的实际使用方式——用户与模型是\textbf{对话}，而非做选择题。

\subsection{评估方法二：Human Evaluation（人类评估）}

\textbf{原理}：人类评审员在两个匿名模型的回答中选择更好的那个（如 Chatbot Arena）。

\textbf{优点}：
\begin{itemize}
    \item 直接反映人类偏好，是 Post-Training 的终极优化目标
    \item 可精细控制评估维度（如只关注毒性、只关注准确性）
    \item 数据污染风险低
\end{itemize}

\textbf{缺点}：
\begin{itemize}
    \item 极难规模化，成本高、耗时长
    \item 人类评审\textbf{有严重偏见}：更长的回答、更自信的语气会被偏好，即使内容有误
\end{itemize}

\begin{warningbox}{人类偏好与自动化 Benchmark 不相关}
讲者展示了一项关键发现：Chatbot Arena 排名与 MMLU、GSM8K、HumanEval 等自动化 benchmark 的相关性非常低。一个 MMLU 高分的模型可能在人类偏好上表现很差，反之亦然。因此，\textbf{必须同时使用两类评估}方法才能全面了解模型能力。
\end{warningbox}

\subsection{评估方法三：LLM-as-Judge}

为了在规模化和人类偏好之间找到折中，可以使用 LLM 替代人类进行评估：

\begin{itemize}
    \item \textbf{单 Judge}：一个强力 LLM（如 GPT-4）对回答评分
    \item \textbf{多 Judge（Jury）}：多个 LLM 组成评审团，投票决定优劣——更稳健可靠
\end{itemize}

\textbf{优点}：与人类偏好高度相关，能处理自由形式的复杂任务，提供直接反馈。

\textbf{缺点}：LLM judge 自身也有偏见（且与人类偏见高度相似），仍需少量人类评估来校验 judge 的一致性。

\subsection{自建评估体系的建议}

\begin{enumerate}
    \item \textbf{尽早开始}——在 fine-tuning 之前就设计评估，类似 test-driven development
    \item \textbf{持续迭代}——评估数据集不是一成不变的，需根据模型表现不断修正
    \item \textbf{混合使用}——自动化 benchmark + LLM-as-Judge + 少量人类评估
    \item \textbf{横向比较}——不仅比较自己的各版本 fine-tune，也与其他模型和架构对比
\end{enumerate}

\begin{knowledgebox}{评估驱动的迭代开发}
讲者强调 Post-Training 的时间分配应大致为：数据准备占三分之一，训练占三分之一，评估占三分之一。评估不是事后验证，而是贯穿整个流程的核心环节。每轮评估的结果直接指导下一轮数据和训练的改进方向。
\end{knowledgebox}

\subsection{本章小结}

LLM 评估主要有三种方法：自动化 benchmark（可扩展但不反映真实使用）、人类评估（最接近目标但难以规模化）、LLM-as-Judge（折中方案）。三者互补，应组合使用。人类偏好与自动化 benchmark 的低相关性是设计评估体系时必须认识到的关键事实。

\newpage
